\documentclass[11pt]{article}
\usepackage[margin=1in]{geometry}
\usepackage{graphicx}
\usepackage{amsmath, amssymb}
\usepackage{booktabs}
\usepackage{listings}
\usepackage{xcolor}
\usepackage{hyperref}
\usepackage{caption}
\usepackage{enumitem}

\definecolor{codebg}{RGB}{245,245,250}
\definecolor{cmdbg}{RGB}{40,42,54}
\definecolor{cmdfg}{RGB}{248,248,242}

\lstdefinestyle{py}{
    language=Python,
    basicstyle=\ttfamily\small,
    backgroundcolor=\color{codebg},
    keywordstyle=\bfseries\color{blue!70!black},
    stringstyle=\color{green!50!black},
    commentstyle=\itshape\color{gray},
    breaklines=true, frame=single, framesep=4pt, showstringspaces=false}

\lstdefinestyle{shell}{
    language=bash,
    basicstyle=\ttfamily\footnotesize\color{cmdfg},
    backgroundcolor=\color{cmdbg},
    breaklines=true, frame=single, framesep=4pt, showstringspaces=false}

\hypersetup{colorlinks=true, urlcolor=blue, linkcolor=black,
            citecolor=black, breaklinks=true}

\title{\Large\bfseries picoacq\\[2pt]
       \large Data Acquisition User Manual\\[8pt]
       \normalsize Coherent multi-channel HF capture with the
       PicoScope 5444D for triloop polarimetry}
\author{}
\date{\today \\[6pt] \small Version 0.1.0}

\begin{document}
\maketitle

\tableofcontents
\bigskip\hrule\bigskip

%-----------------------------------------------------------------------
\section{Overview}
%-----------------------------------------------------------------------

\textsf{picoacq} is the data-acquisition companion to \textsf{triloop},
the analysis package for three-loop magnetic-field-array HF polarimetry.
\textsf{picoacq} drives a Pico Technology PicoScope 5444D (or other
5000-series digitiser) over USB, captures four channels of coherently
sampled IQ-style real signals, and writes them to an HDF5 file format
that \textsf{triloop} reads.

This manual covers acquisition only.  For analysis, polarimetry, and
the underlying magnetoionic physics, see the companion documents
\textsl{triloop User Manual} and \textsl{Three-Loop Magnetic-Field
Array Technical Report}.

\subsection{What \textsf{picoacq} does}

\begin{enumerate}\itemsep=2pt
\item Detects an attached PicoScope 5000-series digitiser via Pico's USB
      driver, opens it at the appropriate vertical resolution
      (Section~\ref{sec:resolution}).
\item Configures the user-selected channels for input range, coupling,
      and sample rate (Sections~\ref{sec:channels}--\ref{sec:rate}).
\item Triggers a block-mode acquisition of the requested duration,
      reads the raw 16-bit ADC samples back to host memory, converts
      them to volts, and writes the result to an HDF5 file with
      complete metadata (Section~\ref{sec:fileformat}).
\item In continuous-monitor mode, repeats this loop on a fixed cadence
      to build a long-running dataset suitable for ionospheric studies.
\end{enumerate}

If no PicoScope is attached or the SDK is missing, \textsf{picoacq}
falls back to a synthetic data simulator so the analysis chain can be
tested without hardware.

\subsection{Audience}

This manual assumes basic Python familiarity, comfort with command-line
tools, and a working knowledge of HF radio signal capture.  The reader
should be familiar with concepts like sample rate, voltage range,
ADC resolution, and the difference between block and streaming
acquisition.

%-----------------------------------------------------------------------
\section{Hardware setup}
%-----------------------------------------------------------------------

\subsection{The PicoScope 5444D}

The 5444D is a 4-channel 200\,MHz USB-connected oscilloscope with a
variable-resolution (``FlexRes'') ADC.  Specifications relevant to
HF capture:

\begin{itemize}\itemsep=2pt
\item \textbf{4 analog input channels}, BNC, switchable between
      $1\,\mathrm{M\Omega}\,\|\,15$\,pF and $50\,\Omega$ termination.
\item \textbf{Variable resolution}, 8 to 16 bits.  See
      Section~\ref{sec:resolution} for the trade-off between resolution
      and channel count.
\item \textbf{Synchronous sampling} of all enabled channels.  All four
      channels share a single sample clock and trigger.  Channel-to-channel
      sample skew is below one sample period.
\item \textbf{Voltage ranges} from $\pm 10$\,mV to $\pm 20$\,V in 11
      steps.
\item \textbf{Streaming throughput} over USB 3.0: $\sim$\,400\,MB/s
      sustained.
\item \textbf{Onboard memory}: 512\,MS total, split across enabled
      channels.  In 4-channel mode each channel gets 128\,MS; in
      $\le 2$-channel mode each gets 256\,MS.  This sets the
      maximum block-mode capture length (8.2\,s for 4 channels at the
      \textsf{picoacq} default rate of 15.625\,MS/s).
\item \textbf{No external reference clock input.}  The 5444D
      enclosure has only the four BNC signal inputs; there is no
      separate 10\,MHz reference jack.  A GPS-disciplined oscillator
      can still be used as a reference by injecting its 10\,MHz pilot
      tone into one (or all) of the input channels through a signal
      combiner --- the recovered pilot then disciplines timing in
      software.
\end{itemize}

\subsection{The crucial connection check}

Before installing any software, plug the PicoScope into the host via
its USB-3 cable and visually verify the LED indicator lights up green
(power) and the host OS sees the device as a USB peripheral.  On macOS:

\begin{lstlisting}[style=shell]
system_profiler SPUSBDataType | grep -i pico
\end{lstlisting}

Should print at least one match line.  If nothing prints, the device is
not detected --- check the USB cable, port (must be USB 3.0+, not 2.0),
and that the LED is illuminated.

\subsection{Sanity-test with the GUI app}

Before driving the scope from \textsf{picoacq}, run Pico's free
GUI application \texttt{PicoScope.app} (download from
\url{https://www.picotech.com/downloads}).  This rules out any hardware
issue independent of \textsf{picoacq} code.

The GUI is also useful for:
\begin{itemize}\itemsep=2pt
\item \textbf{Choosing voltage ranges} per channel by watching live
      time-domain traces from your loop preamps and whip; pick a range
      that uses most of the dynamic range without clipping.
\item \textbf{Checking signal presence} via the live FFT spectrum view,
      e.g.\ confirming that WWV at the expected $\sim$\,2\,kHz IF
      offset is visible above the noise floor.
\item \textbf{Trigger debugging} interactively before scripting.
\end{itemize}

\textbf{Important:} the GUI app and \textsf{picoacq} cannot both
hold the device handle simultaneously.  \textbf{Quit
\texttt{PicoScope.app} before invoking \textsf{picoacq capture}} or
the SDK call will fail with a ``device in use'' error.

%-----------------------------------------------------------------------
\section{Software installation}
\label{sec:install}
%-----------------------------------------------------------------------

\textsf{picoacq} requires three software pieces.  All three are
independent and must be installed separately on macOS, Linux, or
Windows.

\subsection{Layer 1: PicoSDK system library}

The PicoSDK is a system-level shared library installed via Pico's
own installer.  It provides the C-level functions
\texttt{ps5000aOpenUnit}, \texttt{ps5000aSetChannel}, etc.\ that
all higher-level wrappers eventually call.

\begin{enumerate}\itemsep=2pt
\item Visit \url{https://www.picotech.com/downloads}.
\item Select \textbf{PicoSDK} (NOT \textbf{PicoScope}, which is the
      GUI app).
\item Choose your operating system and architecture.
\item Install the resulting \texttt{.pkg} (macOS), \texttt{.deb}
      / \texttt{.rpm} (Linux), or \texttt{.exe} (Windows).
\end{enumerate}

\textbf{Where the library lands on macOS}: as of 2026, Pico packages
the dylibs as a framework at
\texttt{/Library/Frameworks/PicoSDK.framework/Libraries/}.
Each driver lives in its own subdirectory (e.g.\
\texttt{libps5000a/libps5000a.dylib}).  This is the modern macOS
install layout, distinct from the older flat \texttt{/usr/local/lib/}
location.

To verify successful install:
\begin{lstlisting}[style=shell]
sudo find / -name "libps5000a*" 2>/dev/null
\end{lstlisting}

This should list at least one \texttt{libps5000a.dylib} (or platform
equivalent).

\subsection{Layer 2: \texttt{picosdk} Python wrapper}

A thin ctypes wrapper around the system library, distributed via PyPI:

\begin{lstlisting}[style=shell]
python3 -m pip install picosdk
\end{lstlisting}

\textbf{Common pitfall on macOS:} the system has multiple Python
installations (a 3.9 from Apple's macOS, and a 3.13 from
\url{python.org}).  \texttt{pip} and \texttt{pip3} may resolve to
different installations.  Use the \texttt{-m pip} form to force the
binding:

\begin{lstlisting}[style=shell]
python3 -m pip install picosdk
\end{lstlisting}

Verify:
\begin{lstlisting}[style=shell]
python3 -c "from picosdk.ps5000a import ps5000a; print('SDK OK')"
\end{lstlisting}

\subsection{Layer 3: \texttt{picoacq} itself}

Editable install from the project tree:

\begin{lstlisting}[style=shell]
cd <project-root>
pip install -e ./picoacq
\end{lstlisting}

This installs the \texttt{picoacq} command-line tool and Python
package.

\subsection{Verifying the install}

\begin{lstlisting}[style=shell]
picoacq --help
picoacq capture --simulate -o /tmp/test.h5 --duration 0.1
\end{lstlisting}

The first should print a help message; the second should write a
synthetic capture file using the simulator.  No real hardware needed
for either.

\textsf{picoacq} automatically detects the macOS Framework path and
adds it to \texttt{DYLD\_LIBRARY\_PATH} at import time, so the
common ``library not found'' error on macOS is handled transparently.
On Linux you may need to add the SDK's library directory to
\texttt{LD\_LIBRARY\_PATH} manually if the installer didn't.

%-----------------------------------------------------------------------
\section{Hardware capabilities and resolution selection}
\label{sec:resolution}
%-----------------------------------------------------------------------

The 5444D's FlexRes ADC trades vertical resolution against the number
of simultaneously-enabled channels and the per-channel sample rate.
Knowing which mode you're in is essential to avoid the
``\texttt{PICO\_INVALID\_NUMBER\_CHANNELS\_FOR\_RESOLUTION}'' error
when configuring your acquisition.

\begin{table}[h]
\centering
\begin{tabular}{lccc}
\toprule
\textbf{Resolution} & \textbf{Max channels} & \textbf{Max rate per ch} & \textbf{Use for} \\
\midrule
8-bit & 4 & 1\,GS/s (1 ch) / 250\,MS/s (4 ch) & High bandwidth, low precision \\
12-bit & 4 & 500\,MS/s & General oscilloscope use \\
14-bit & 4 & 125\,MS/s & \textbf{HF triloop polarimetry} \\
15-bit & 2 & 125\,MS/s & 2-ch high-precision \\
16-bit & 2 & 62.5\,MS/s & Ultra-precision, 2-ch only \\
\bottomrule
\end{tabular}
\caption{PicoScope 5444D resolution / channel-count compatibility.
Note: 16-bit mode is restricted to 2 channels.  For
4-channel triloop applications, 14-bit is the highest available
resolution.}
\label{tab:resolution}
\end{table}

\textbf{For triloop polarimetry with 3 loops + 1 whip = 4 channels,
\textsf{picoacq} automatically selects 14-bit resolution.}  This gives
you 86\,dB of dynamic range, which is far more than necessary for HF
work after narrow-band filtering.  The 16-bit mode that some Pico
documentation references is a 2-channel-only mode and is not used by
\textsf{picoacq}.

If you have only 2 (or fewer) channels enabled, \textsf{picoacq}
auto-selects 16-bit resolution.  This selection happens in
\texttt{picoacq.ps5444a.open\_5444()} based on the channel count
passed by the recorder.

%-----------------------------------------------------------------------
\section{Channel configuration}
\label{sec:channels}
%-----------------------------------------------------------------------

Each enabled channel needs three settings: a \textbf{voltage range}
(per-channel, $\pm 10$\,mV to $\pm 20$\,V), a \textbf{coupling mode}
(AC or DC), and an \textbf{enable bit} (\textsf{picoacq} sets this
automatically for any channel listed in \texttt{--channels}).

\subsection{Voltage range selection}

The voltage range determines the ADC's full-scale span.  The 14-bit
ADC has $2^{14} = 16384$ codes; for a $\pm 2$\,V range the code-to-volts
factor is $4 / 16384 \approx 244\,\mu$V/code.

\textbf{Rule of thumb}: pick a range so the largest expected signal
peak uses about 50--80\% of full scale.  Too small a range clips the
signal and produces nonlinearities; too large a range wastes ADC
bits on noise.

For HF triloop work with active mag-loops and a vertical mini-whip,
typical signal levels at the receiver after amplification are 50\,mV
to 1\,V peak.  $\pm 2$\,V is usually a safe default.

\subsection{Auto-ranging}
\label{sec:autorange}

By default \textsf{picoacq} performs an \textbf{auto-range probe}
before each capture, eliminating the need to know your signal levels
in advance.  The procedure:

\begin{enumerate}\itemsep=2pt
\item Open the scope at the maximum range ($\pm 20$\,V) on every
      channel.
\item Take a brief probe capture (default 100\,ms at the requested
      sample rate).
\item For each channel: measure the peak $|x|$ and the RMS of the
      captured samples in volts.
\item For each channel: pick the smallest available Pico voltage range
      whose full scale $\geq K \cdot \mathrm{peak}$ (default
      $K = 3$, i.e.\ 10\,dB headroom on top of the measured peak).
\item Re-configure each channel at the chosen range and run the
      science capture.
\end{enumerate}

The Pico's available ranges are quantised:
$\pm 10$\,mV, $\pm 20$, $\pm 50$, $\pm 100$, $\pm 200$, $\pm 500$\,mV,
$\pm 1$, $\pm 2$, $\pm 5$, $\pm 10$, $\pm 20$\,V.  The auto-range
chooser always rounds \emph{up} to the next available range, so the
actual headroom may exceed $K$ by up to a factor of $\sim 2$ depending
on where your peak falls between range steps.

\paragraph{What auto-ranging buys you.}  At the maximum range
($\pm 20$\,V), a 50\,mV peak signal occupies only $0.25\%$ of full
scale.  With a 14-bit ADC this means only $\sim 5$ effective bits of
useful resolution = $\sim 32$\,dB quantization-limited SNR.
Auto-ranging this same signal to (say) $\pm 200$\,mV yields a 5×
headroom and gives you the full 14 bits across the active range:
$\sim 78$\,dB SNR.  That's a $\sim 46$\,dB improvement in dynamic
range from one decision.

\paragraph{Disabling auto-range.}  If you want to set ranges manually
(e.g., when calibrating against a known signal source where you
control the input amplitude), pass \texttt{--no-auto-range} to the
CLI:

\begin{lstlisting}[style=shell]
# Manual ranges (--range applies to all enabled channels):
picoacq capture -o cap.h5 --duration 8 \
    --channels A,B,C,D --range 0.5 --no-auto-range
\end{lstlisting}

\paragraph{Tuning the headroom factor.}  The default $K = 3$ is
moderately conservative (10\,dB above peak).  For very stable signals
where transient excursions are unlikely, you can try $K = 2$ to gain
$\sim 3$\,dB more dynamic range.  For unpredictable signals (e.g.,
picking up sferics from nearby thunderstorms),
$K = 5$ to $10$ is safer.  Override with \texttt{--headroom}:

\begin{lstlisting}[style=shell]
picoacq capture --headroom 5 ...
\end{lstlisting}

\paragraph{Probe-result metadata.}  The HDF5 file's
\texttt{capture\_settings} dict includes:
\begin{itemize}\itemsep=2pt
\item \texttt{auto\_range\_enabled}: True / False
\item \texttt{headroom}: the $K$ factor used
\item \texttt{auto\_range\_probe}: the per-channel \texttt{peaks\_volts}
      and \texttt{rms\_volts} measured during the probe
\item \texttt{picoscope\_ranges\_volts}: what the Pico's range was
      actually set to (after rounding up to the available range)
\end{itemize}
This lets you reconstruct the dynamic range used after the fact.

\subsection{Coupling}

DC coupling preserves the full signal including any DC offset.  AC
coupling adds a series capacitor that high-passes at $\sim 1$\,Hz.

For HF capture, DC and AC are both fine.  AC removes any DC offset
from the active electronics, which is useful for keeping the ADC
centered.  Default in \textsf{picoacq} is DC.

\subsection{Channel-name to PicoScope-port mapping}

The PicoScope 5444D has BNC connectors labelled \texttt{A},
\texttt{B}, \texttt{C}, \texttt{D}.  By convention in \textsf{triloop}:
\begin{itemize}\itemsep=2pt
\item \texttt{A}: loop L1 (north-of-up)
\item \texttt{B}: loop L2 ($+120^\circ$)
\item \texttt{C}: loop L3 ($-120^\circ$)
\item \texttt{D}: vertical whip (or other reference)
\end{itemize}
You may rearrange via \texttt{loops\_config} in your analysis script.

%-----------------------------------------------------------------------
\section{Sample-rate selection}
\label{sec:rate}
%-----------------------------------------------------------------------

The PicoScope 5000-series uses an integer ``timebase'' index to specify
sample period.  At 14-bit resolution with 4 channels active, the
available sample periods are quantized to specific values; you
\emph{request} a target rate, and the SDK returns the closest available.

\subsection{Default rate: 15.625\,MS/s, all WWV bands at once}
\label{sec:defaultrate}

\textsf{picoacq}'s default sample rate is \textbf{15.625\,MS/s} in
4-channel mode.  This rate was chosen because it is high enough that
\emph{all five WWV bands} (2.5, 5, 10, 15, and 20\,MHz) fold to
non-overlapping baseband locations through bandpass undersampling, so
a single capture digitises every band simultaneously and coherently
across all four input channels.

Under bandpass undersampling, an RF tone at $f_\text{rf}$ aliases
to baseband as
\begin{equation*}
  f_\text{bb} \;=\; \bigl|\, f_\text{rf} - n \cdot f_s \,\bigr|,
  \qquad n = \mathrm{round}(f_\text{rf} / f_s).
\end{equation*}
The Nyquist zone $z = \lfloor 2 f_\text{rf}/f_s \rfloor + 1$
determines whether the spectrum is preserved (odd zones) or
frequency-inverted (even zones).  At $f_s = 15.625$\,MS/s the WWV
bands map as follows:

\begin{table}[h]
\centering
\begin{tabular}{cccc}
\toprule
\textbf{RF (MHz)} & \textbf{Nyquist zone} & \textbf{Baseband (MHz)} & \textbf{Inverted?} \\
\midrule
 2.5  & 1 & 2.500 & no  \\
 5.0  & 1 & 5.000 & no  \\
10.0  & 2 & 5.625 & yes \\
15.0  & 2 & 0.625 & yes \\
20.0  & 3 & 4.375 & no  \\
\bottomrule
\end{tabular}
\caption{WWV bands at $f_s = 15.625$\,MS/s.  All five bands occupy
distinct, non-overlapping baseband locations.  Bands in even Nyquist
zones come out spectrum-inverted, which the analysis tools handle
via complex conjugation of the appropriate down-converted slice.}
\label{tab:alias}
\end{table}

The helper \texttt{picoacq.alias.alias\_of(f\_rf, f\_s)} computes
this mapping programmatically; \texttt{rf\_for\_baseband()} performs
the inverse lookup given a list of expected RF bands (typically the
\texttt{rf\_bands\_hz} field stored in the file's
\texttt{capture\_settings}).

\paragraph{Why not 12.5\,MS/s?}  An apparently-tempting lower rate of
12.5\,MS/s has the WWV 5 and 10\,MHz bands both folding to 2.5\,MHz
baseband on top of the unaliased 2.5\,MHz signal --- three bands pile
up in the same baseband bin and become inseparable.  15.625\,MS/s is
the lowest rate at which the five WWV bands separate cleanly.

\paragraph{Why not 62.5\,MS/s direct?}  Sampling above $2 \cdot 20$\,MHz
puts every WWV band in zone 1 (no aliasing), but the 128\,MS/channel
buffer then holds only $\sim 2$\,s of data.  15.625\,MS/s preserves
absolute phase information at all five bands while extending the
maximum block to 8.2\,s --- the longer block dominates phase-stability
science requirements ($\tau \le 1$\,s coherence times).

\subsection{The \texttt{--rf-bands} flag}

By default, \texttt{picoacq capture} and \texttt{picoacq monitor}
record the WWV band list
\texttt{2.5e6,5e6,10e6,15e6,20e6} into the file's
\texttt{capture\_settings} as \texttt{rf\_bands\_hz}.  Override with
\texttt{--rf-bands} for a different observation plan
(e.g.\ CHU at 3.33, 7.85, 14.67\,MHz):

\begin{lstlisting}[style=shell]
picoacq capture --rf-bands "3.33e6,7.85e6,14.67e6" ...
\end{lstlisting}

Pass an empty string to omit the metadata entirely.  Downstream tools
read this field to map detected baseband peaks back to the originating
RF band.

\subsection{Common sample rates and their timebase indices}

For 4-channel 14-bit operation:

\begin{table}[h]
\centering
\begin{tabular}{cccc}
\toprule
\textbf{Timebase} & \textbf{Sample period} & \textbf{Sample rate} & \textbf{Useful for} \\
\midrule
3 & 64\,ns & 15.625\,MS/s & \textbf{picoacq default}; all WWV bands at once \\
4 & 96\,ns & 10.42\,MS/s & --- \\
8 & 128\,ns & 7.81\,MS/s & --- \\
512 & 4.08\,$\mu$s & 245\,kHz & Single-band IF capture \\
1024 & 8.18\,$\mu$s & 122\,kHz & Long captures, less data \\
4000 & 31.98\,$\mu$s & 31.3\,kHz & Audio-rate analysis \\
\bottomrule
\end{tabular}
\caption{Common timebase indices for 4-channel 14-bit
acquisition on the 5444D.  Timebase indices below 3 are invalid
in 4-channel mode and are skipped.  See the SDK reference for the
full table.}
\label{tab:timebase}
\end{table}

\textsf{picoacq capture --rate <Hz>} accepts any positive value and
selects the timebase whose actual rate is at or below your request.
The rate \emph{actually used} is recorded in the HDF5 metadata as
\texttt{sample\_rate}.  Don't assume your requested rate equals the
actual rate; always check the file's metadata.

\subsection{Memory budget and block duration limits}

Samples are stored as float32 (4 bytes each) in volts.  A
\textbf{block-mode capture cannot exceed the per-channel onboard
memory}; \textsf{picoacq} validates the requested duration against
the buffer up front and refuses with a clear error if you ask for
too much:

\begin{table}[h]
\centering
\begin{tabular}{ccc}
\toprule
\textbf{Channels} & \textbf{Buffer per ch} & \textbf{Max block @ 15.625\,MS/s} \\
\midrule
4 & 128\,MS & 8.19\,s \\
2 & 256\,MS & 16.38\,s \\
1 & 512\,MS & 32.77\,s \\
\bottomrule
\end{tabular}
\caption{Onboard sample budget per channel and the resulting maximum
block duration at the default rate.}
\label{tab:bufferbudget}
\end{table}

A typical 8\,s capture at 15.625\,MS/s with 4 channels stores:
\[
8 \cdot 15{,}625{,}000 \cdot 4 \cdot 4 \;=\; 2{,}000\,\text{MB}
\]
on disk before compression.  HDF5's gzip cuts this $\sim$30\% for
typical signals, but at this volume host RAM is the limiting factor.
For long captures or higher rates, switch to streaming mode (not yet
implemented in \textsf{picoacq}; planned for v0.2).



%-----------------------------------------------------------------------
\section{Command-line interface}
%-----------------------------------------------------------------------

\textsf{picoacq} ships a single CLI tool with two main subcommands:
\texttt{capture} and \texttt{monitor}.  Both produce HDF5 capture files
in the format described in Section~\ref{sec:fileformat}.

\subsection{One-shot capture: \texttt{picoacq capture}}

\begin{lstlisting}[style=shell]
picoacq capture -o cap.h5 \
    --duration 8 \                    # seconds (max 8.19 at 4ch/15.625 MS/s)
    --rate     15625000 \             # samples per second (default)
    --channels A,B,C,D \              # comma-separated, no spaces
    --range    2.0 \                  # +/- volts (used only if --no-auto-range)
    --coupling DC                     # or AC
\end{lstlisting}

All flags above are at their defaults except \texttt{-o}; bare
\texttt{picoacq capture} produces an 8-second 4-channel
15.625\,MS/s capture written to the cwd.

\paragraph{Output path defaults.}  If \texttt{--output} (or \texttt{-o})
is omitted, the file is written into the \textbf{current working
directory} as \texttt{capture\_<UTC-timestamp>.h5}.  This makes
\texttt{cd} into a directory on an external disk and running
\texttt{picoacq capture} the natural way to deposit data there:

\begin{lstlisting}[style=shell]
cd /Volumes/MyExternalDisk/triloop_data/
picoacq capture --duration 8 --rate 15625000 \
    --channels A,B,C,D
# -> writes /Volumes/MyExternalDisk/triloop_data/capture_20260523T173951Z.h5
\end{lstlisting}

If you pass an explicit \texttt{-o cap.h5}, that's interpreted relative
to your current directory.  Pass an absolute path
(\texttt{-o /Volumes/Disk/cap.h5}) to override the cwd-relative
default.

If the PicoSDK or hardware is unavailable, \textsf{picoacq} prints a
warning and falls back to its simulator (Section~\ref{sec:simulator}).
A successful real-hardware capture writes only the file path to stdout
and creates the HDF5 file.

\subsection{Continuous monitoring: \texttt{picoacq monitor}}

For long unattended runs, \texttt{picoacq monitor} takes one capture
every \texttt{--interval} seconds and writes a timestamped file:

\begin{lstlisting}[style=shell]
picoacq monitor -d /data/triloop_runs/ \
    --interval 600 \                  # seconds between capture starts
    --duration 8 \                    # length of each capture
    --rate     15625000 \             # all WWV bands at once
    --prefix   wwv_                   # filename prefix
\end{lstlisting}

Output filenames look like
\texttt{wwv\_20260520T144312Z.h5}.

\paragraph{Output directory defaults.}  Like \texttt{capture}, the
\texttt{--out-dir} option for \texttt{monitor} defaults to the current
working directory (.).  This means a typical workflow is:

\begin{lstlisting}[style=shell]
cd /Volumes/ExternalDisk/wwv_run_2026_05/
picoacq monitor --interval 600 --duration 8 --rate 15625000 \
    --channels A,B,C,D --prefix wwv_
# -> writes wwv_<timestamp>.h5 files into the external disk
\end{lstlisting}

Pass \texttt{-d /some/other/path} to override.  SIGTERM and SIGINT are handled
cleanly --- the current capture finishes and the loop exits.  Pass
\texttt{--n-captures N} to stop after $N$ captures rather than running
indefinitely.

\subsection{Simulator mode: \texttt{picoacq capture --simulate}}
\label{sec:simulator}

The simulator generates a synthetic three-loop dataset with controllable
arrival direction, polarization, additive Gaussian noise, and Faraday
rotation rate.  It writes the same HDF5 format as the real-hardware
path so \textsf{triloop} can analyze the result without modification.

\begin{lstlisting}[style=shell]
picoacq capture -o cap_sim.h5 --simulate \
    --sim-pol     rcp \                  # linear_vertical, linear_horizontal,
                                          # rcp, lcp, elliptical
    --sim-faraday 30 \                   # deg/s Faraday rotation rate
    --sim-snr     25 \                   # per-channel SNR (dB)
    --sim-az      273 \                  # arrival azimuth (deg from N CW)
    --sim-el      12 \                   # arrival elevation (deg)
    --sim-carrier 25000                  # IF carrier freq (Hz)
\end{lstlisting}

The simulator is also used as the no-hardware fallback in real-capture
mode.  This is useful for code development on a laptop without a scope
attached.

%-----------------------------------------------------------------------
\section{HDF5 file format}
\label{sec:fileformat}
%-----------------------------------------------------------------------

\textsf{picoacq} writes a self-describing HDF5 file (format version
0.1) that is also the input format for \textsf{triloop}.

\begin{lstlisting}[style=shell]
/                            HDF5 root
  @format_version            "0.1"
/metadata                    group, with the following attributes:
  @sample_rate               float64, Hz (actual achieved rate)
  @start_time_utc            ISO-8601 string, e.g. "2026-05-23T17:33:37.994Z"
  @duration_s                float64
  @scope_model               string, e.g. "PicoScope 5444D"
  @scope_serial              string
  @capture_settings          JSON-encoded dict (full kwargs)
  @loops_config              JSON-encoded dict (loop geometry)
/channels                    group containing per-channel datasets
  /A                         float32 array, length N (in volts)
  /B                         (same)
  /C                         (same)
  /D                         (optional reference channel)
  Each channel dataset has attributes:
    @gain_v_per_count        float64  (1.0; values are already in volts)
    @offset_v                float64  (typically 0.0)
/time                        optional float64 dataset, length N
                              (regenerated from sample_rate if absent)
\end{lstlisting}

The format is intentionally simple --- a third party can write or read
it with $\sim$\,30 lines of Python using \texttt{h5py}.

\subsection{Reading a capture file in Python}

\begin{lstlisting}[style=py]
import h5py, json, numpy as np

with h5py.File("cap.h5", "r") as f:
    sr = float(f["metadata"].attrs["sample_rate"])
    duration = float(f["metadata"].attrs["duration_s"])
    channels = {name: f["channels"][name][:].astype(np.float64)
                for name in f["channels"].keys()}
    loops_config = json.loads(f["metadata"].attrs["loops_config"])

t = np.arange(len(next(iter(channels.values())))) / sr
print(f"Loaded {len(channels)} channels, {len(t)} samples each, "
      f"at {sr:.1f} Hz over {duration:.2f} s")
\end{lstlisting}

For more complex use, \texttt{triloop.io\_hdf5.read\_capture()}
provides a tested implementation that returns a dict with all metadata
already parsed.

%-----------------------------------------------------------------------
\section{Worked examples}
%-----------------------------------------------------------------------

\subsection{Example 1: smoke-test on the bench}

After installation, with the PicoScope plugged in via USB and
\texttt{PicoScope.app} not running:

\begin{lstlisting}[style=shell]
# 0.5-second snapshot, all 4 channels, auto-range (default)
picoacq capture --duration 0.5 --channels A,B,C,D

# Inspect the result
triloop summary capture_*.h5
\end{lstlisting}

The auto-range probe runs first, then the science capture:
\begin{lstlisting}[style=shell]
[picoacq] auto-range probe (100 ms at 15625000 Hz, headroom=3.0)
[picoacq] probe results:
           A: peak=230.11 mV  RMS=107.27 mV  -> range +/- 1.0 V
           B: peak=234.99 mV  RMS=113.19 mV  -> range +/- 1.0 V
           C: peak=272.23 mV  RMS=126.72 mV  -> range +/- 1.0 V
           D: peak=366.83 mV  RMS=183.37 mV  -> range +/- 2.0 V
wrote /Users/.../capture_20260523T173951Z.h5
\end{lstlisting}

\texttt{triloop summary} confirms:
\begin{lstlisting}[style=shell]
file:           capture_20260523T173951Z.h5
format_version: 0.1
start_time_utc: 2026-05-23T17:33:37.994093+00:00
sample_rate:    209732 Hz
duration_s:     0.4768 s
scope:          PicoScope 5444D (s/n EW123/4567)
channels:
  A: n=100000, range=[-0.005, +0.023]
  B: n=100000, range=[-0.003, +0.024]
  C: n=100000, range=[-0.003, +0.027]
  D: n=100000, range=[-0.001, +0.018]
\end{lstlisting}

The voltage ranges shown by \texttt{summary} are the actual
\textbf{captured-signal} range during the science capture, which is
typically much smaller than the auto-ranged Pico setting (because the
auto-range probe captured a different instant of input transients,
and we still rounded up to a quantized Pico range).

The actual achieved sample rate may differ slightly from the
requested rate --- the SDK selects the nearest valid timebase index
at or below the request (Section~\ref{sec:rate}).  Always consult
\texttt{metadata/sample\_rate} in the file rather than the requested
value when computing time axes.

\subsection{Example 2: 8-second all-WWV-band capture}

Connect a triloop array (3 mag-loops to channels A-C) plus a vertical
whip (channel D), with each input from an active preamp.  Take an
8-second capture at the default 15.625\,MS/s, which simultaneously
digitises all five WWV bands (Section~\ref{sec:defaultrate}):

\begin{lstlisting}[style=shell]
picoacq capture -o /data/triloop/cap_wwv_$(date -u +%Y%m%dT%H%M%SZ).h5 \
    --duration 8 --channels A,B,C,D
\end{lstlisting}

Resulting file size: about 1.4\,GB compressed (raw is 2.0\,GB).
With $\sim$125\,M samples per channel over 8 seconds, the WWV bands
appear at the baseband locations listed in Table~\ref{tab:alias}.

Run \textsf{triloop} analysis on it:

\begin{lstlisting}[style=shell]
triloop analyze cap_wwv_20260523T173337Z.h5 \
    --rf-band 5e6 \        # one of the recorded rf_bands_hz entries
    --bw 4000 \
    --az 9 --el 35         # FC->APO geometry
\end{lstlisting}

\textsf{triloop} reads the file's \texttt{rf\_bands\_hz} setting,
maps the requested RF band through \texttt{picoacq.alias} to its
baseband position, and slices that band out for analysis.

\subsection{Example 3: continuous overnight monitor at APO}

Run continuously, 8 seconds out of every 10 minutes:

\begin{lstlisting}[style=shell]
picoacq monitor -d /mnt/data/wwv_monitor/ \
    --interval 600 --duration 8 --channels A,B,C,D \
    --prefix wwv_
\end{lstlisting}

Files appear as \texttt{wwv\_20260523T173337Z.h5},
\texttt{wwv\_20260523T174337Z.h5}, etc.  At $\sim$1.4\,GB per file
this is about 8.5\,GB/hour, or $\sim$\,200\,GB per 24-hour
observation --- pre-allocate disk accordingly.

The monitor catches SIGINT and SIGTERM cleanly, finishing the current
capture before exiting.  Use \texttt{Ctrl-C} or \texttt{kill -TERM PID}
to stop.

%-----------------------------------------------------------------------
\section{Troubleshooting}
%-----------------------------------------------------------------------

\subsection{Common errors and what they mean}

\paragraph{\texttt{[picoacq] real capture failed: picosdk not installed}}
The \texttt{picosdk} Python wrapper is missing.  Run
\texttt{python3 -m pip install picosdk}.  Note the use of
\texttt{python3 -m pip} rather than \texttt{pip} --- on macOS the
two may resolve to different Python installations.

\paragraph{\texttt{[picoacq] real capture failed: PicoSDK (ps5000a) not found, check LD\_LIBRARY\_PATH}}
The system-level PicoSDK library isn't installed, or
\textsf{picoacq}'s framework-path detection didn't find it.  Run
\texttt{sudo find / -name "libps5000a*" 2>/dev/null} to confirm the
library is present on the system.  If it's there but \textsf{picoacq}
isn't finding it, manually set the path:
\begin{lstlisting}[style=shell]
export DYLD_LIBRARY_PATH=/Library/Frameworks/PicoSDK.framework/Libraries/libps5000a:$DYLD_LIBRARY_PATH
\end{lstlisting}

\paragraph{\texttt{PICO\_NOT\_FOUND}}
The PicoScope hardware isn't visible to the SDK.  Possible causes:
\begin{itemize}\itemsep=2pt
\item PicoScope.app or another program is holding the device handle.
      Quit those.
\item USB cable is bad or unplugged.
\item Pico's user-space USB plugin wasn't authorised on macOS.  Re-run
      the installer or check System Preferences $\to$ Security and
      authorise the Pico kext/extension.
\end{itemize}

\paragraph{\texttt{PICO\_INVALID\_NUMBER\_CHANNELS\_FOR\_RESOLUTION}}
You're trying to use 4 channels at 16-bit resolution.  The 5444D
supports at most 2 channels at 16-bit; for 4 channels, you need
14-bit.  \textsf{picoacq} normally auto-selects this; if you've
manually overridden the resolution, choose 14-bit.

\paragraph{\texttt{argument N: TypeError: '\_ctypes.CArgObject' object cannot be interpreted as an integer}}
A ctypes argument-type mismatch in the picosdk Python wrapper.
Should not occur with current \textsf{picoacq}; if you see it, it
likely means the picosdk version on your system has function
signatures that differ from what \textsf{picoacq} expects.  File an
issue with the SDK version (\texttt{python3 -c "import picosdk; print(picosdk.\_\_version\_\_)"}).

\paragraph{\texttt{could not find a timebase for X Hz}}
The requested sample rate isn't achievable with the current channel
count and resolution.  At 4-channel 14-bit, valid sample rates start
around 31\,Hz (timebase 4000) and go up to 125\,MS/s (timebase 3).
If you really need a sample rate \textsf{picoacq} reports as
unreachable, try reducing the channel count or check the
PicoScope 5000-series Programmer's Guide for valid timebase indices.

\subsection{Verifying timebase availability manually}

To diagnose timebase issues, this snippet walks the available
indices for your current channel count:

\begin{lstlisting}[style=py]
import ctypes
from picoacq.ps5444a import open_5444, configure_channel
from picosdk.constants import PICO_STATUS_LOOKUP

with open_5444(n_channels=4) as (ps, handle):
    for ch in ['A','B','C','D']:
        configure_channel(ps, handle, ch, range_volts=2.0)
    print('all 4 channels configured')
    for tb in [3, 4, 8, 16, 64, 256, 512, 1024, 4000]:
        timeIntervalNs = ctypes.c_float()
        maxSamples = ctypes.c_int32()
        ret = ps.ps5000aGetTimebase2(handle,
            ctypes.c_uint32(tb), ctypes.c_int32(1024),
            ctypes.byref(timeIntervalNs),
            ctypes.byref(maxSamples),
            ctypes.c_uint32(0))
        if ret == 0:
            rate = 1.0 / (timeIntervalNs.value * 1e-9)
            print(f"  tb={tb}: dt={timeIntervalNs.value:.3f} ns, "
                  f"rate={rate:.0f} Hz")
        else:
            print(f"  tb={tb}: status={ret} ({PICO_STATUS_LOOKUP[ret]})")
\end{lstlisting}

%-----------------------------------------------------------------------
\section{Hardware capture process, step by step}
%-----------------------------------------------------------------------

For reference, here's the sequence of SDK calls \textsf{picoacq}
makes to perform one capture.  This is the part of \texttt{picoacq.ps5444a}
that ultimately runs:

\begin{enumerate}\itemsep=2pt
\item \texttt{ps5000aOpenUnit}: opens the device handle.  Resolution
      argument is set based on channel count (14-bit for 4-channel
      mode, 16-bit for $\le 2$).
\item For each enabled channel: \texttt{ps5000aSetChannel} configures
      the channel for the requested voltage range and coupling
      (DC/AC).
\item \texttt{ps5000aGetTimebase2}: walks timebase indices and finds
      the one whose sample rate is at or just below the requested
      target rate.  Some early indices are invalid in 4-channel mode
      and are skipped.
\item For each enabled channel: \texttt{ps5000aSetDataBuffer} allocates
      a host-side buffer and tells the driver where to write samples.
\item \texttt{ps5000aRunBlock}: triggers the block acquisition.  The
      driver fills the buffers and signals completion.
\item Poll \texttt{ps5000aIsReady} every 5\,ms until the driver
      reports the block is complete.
\item \texttt{ps5000aGetValues}: pulls the captured samples into the
      pre-allocated host buffers.
\item Convert int16 ADC codes to volts using the per-channel range
      setting (gain = range\_volts / 32767).
\item \texttt{ps5000aCloseUnit}: releases the device handle.
\end{enumerate}

The total time for a 1-second capture at 200\,kSPS is about 1.5
seconds (including 0.5\,s of overhead for opening, configuring, and
closing).  For longer captures the overhead is small.

%-----------------------------------------------------------------------
\section{Limitations of v0.1}
%-----------------------------------------------------------------------

\begin{itemize}\itemsep=2pt
\item \textbf{Block mode only}.  Continuous streaming over USB at
      hundreds of MB/s is not yet implemented.  For now, captures are
      finite-duration block mode reads.  This is fine for typical
      monitor cadences (1\,minute / 10\,minutes) but not for
      uninterrupted multi-hour streaming.
\item \textbf{Same manual-range value across all channels} (when
      \texttt{--no-auto-range} is used).  The CLI takes a single
      \texttt{--range} value applied to every channel.  Auto-range
      always picks per-channel ranges independently.  Per-channel
      manual ranges are supported in the underlying
      \texttt{configure\_channel()} function, just not exposed in the
      CLI.
\item \textbf{No external trigger configuration}.  The current code
      uses a software trigger immediately on \texttt{RunBlock}.
      External hardware-trigger support is on the roadmap for v0.2.
\item \textbf{Serial number reads as a placeholder}.  The metadata
      records ``\texttt{(read from scope handle)}'' rather than the
      actual unit serial.  We have the SDK call needed
      (\texttt{ps5000aGetUnitInfo}); just not wired up yet.
\item \textbf{No software pilot-tone discipline yet}.  The 5444D has
      no external 10\,MHz reference jack; the planned workflow is to
      inject a GPSDO 10\,MHz pilot into one of the input channels and
      correct the recorded signals' time base in software at analysis
      time.  That correction is not yet implemented in
      \textsf{picoacq} or \textsf{triloop}.
\end{itemize}

These will be addressed in subsequent releases.  For most triloop
science applications, the v0.1 feature set is sufficient.

%-----------------------------------------------------------------------
\section{Quick reference card}
%-----------------------------------------------------------------------

\begin{lstlisting}[style=shell]
# Install (one time)
python3 -m pip install picosdk
pip install -e <project>/picoacq

# Quick smoke test
picoacq capture --simulate -o /tmp/test.h5 --duration 0.5

# Single capture (auto-range by default; all WWV bands; written to cwd)
picoacq capture                         # 8 s, 4 ch, 15.625 MS/s

# Single capture, manual range, explicit output
picoacq capture -o cap.h5 --duration 8 \
    --channels A,B,C,D --no-auto-range --range 2.0

# Continuous monitor (auto-range each capture, written to cwd)
picoacq monitor --interval 600 --duration 8 --channels A,B,C,D

# CHU bands instead of WWV (override the default --rf-bands list)
picoacq capture --rf-bands "3.33e6,7.85e6,14.67e6"

# Stop monitor cleanly: Ctrl-C
\end{lstlisting}

\begin{lstlisting}[style=py]
# Read a file in Python
import h5py
with h5py.File("cap.h5", "r") as f:
    sr = float(f["metadata"].attrs["sample_rate"])
    A  = f["channels"]["A"][:]
\end{lstlisting}

\end{document}
